\documentclass[../../../include/open-logic-section]{subfiles}

\begin{document}

\olfileid{sth}{story}{blv}

\olsection{परिशिष्ट: फ्रेगे यांचा मूलनियम पाच}

\olref[rus]{sec} मध्ये आपण पाहिले की रसेल यांनी आपली विरोधापत्ती फ्रेगे यांनी \emph{Grundgesetze} मध्ये मांडलेल्या प्रणालीसमोरील अडचण म्हणून सांगितली. फ्रेगे यांच्या प्रणालीत अनौपचारिक गुणधर्माधारित संचरचनेची थेट मांडणी नव्हती. म्हणून या परिशिष्टात त्या प्रणालीत \emph{नेमके काय} होते, त्याचा अनौपचारिक गुणधर्माधारित संचरचनेशी कसा संबंध आहे आणि रसेल यांच्या विरोधापत्तीशी कसा संबंध आहे, हे अगदी थोडक्यात पाहू.

फ्रेगे यांची प्रणाली \emph{द्वितीय-क्रमाची} आहे. तिचा हेतू \emph{संकल्पनेचा विस्तार} ही संकल्पना औपचारिकरीत्या मांडणे हा होता.\footnote{काटेकोरपणे सांगायचे तर फ्रेगे यांनी अधिक व्यापक संकल्पना, म्हणजे फलनाची ``मूल्यव्याप्ती'', औपचारिकरीत्या मांडण्याचा प्रयत्न केला. संकल्पनांचे विस्तार हे त्या व्यापक संकल्पनेचे विशेष उदाहरण आहे. तपशीलांसाठी \citet[pp.\ 8--9]{Heck2012} पाहा.} फ्रेगे यांच्या संकेतपद्धतीच्या धर्तीवर $\fregeext{x}{F(x)}$ हे \emph{संकल्पना~$F$ चा विस्तार} दर्शवेल. हे साधन ``$F$'' या \emph{विधेयाचे} (प्रथम-क्रमातील) ``$\fregeext{x}{F(x)}$'' या \emph{पदात} रूपांतर करते. या साधनाच्या आधारे फ्रेगे यांनी सदस्यत्वाची पुढील \emph{व्याख्या} दिली:
\[
	a \in b =_\text{df} \exists G(b = \fregeext{x}{G(x)} \land Ga)
\]
स्थूलमानाने, $a \in b$ तेव्हाच आणि तेव्हाच जेव्हा $a$ अशी एखादी संकल्पना पूर्ण करते की तिचा विस्तार $b$ आहे. (येथे ``$\exists G$'' हा संख्यापक द्वितीय-क्रमाचा आहे, हे लक्षात घ्या.) फ्रेगे यांनी \emph{मूलनियम पाच} म्हणून ओळखले जाणारे पुढील तत्त्वही मानले:
$$\fregeext{x}{F(x)} = \fregeext{x}{G(x)} \liff \forall x (Fx \liff Gx)$$
स्थूलमानाने, दोन संकल्पनांचे विस्तार एकच असतात तेव्हाच आणि तेव्हाच त्या नेमक्या त्याच वस्तूंना लागू पडतात. (पुन्हा, ``$F$'' आणि ``$G$'' ही दोन्ही विधेयस्थानी आहेत.) आता एक साधे तत्त्व सदस्यत्वाचा संबंध गुणधर्मपूर्तीशी जोडते:

\begin{lem}[\emph{Grundgesetze} मध्ये]\ollabel{lem:Fregeextensions}
$\forall F \forall a(a \in \fregeext{x}{F(x)} \liff Fa)$
\end{lem}

\begin{proof}
$F$ आणि $a$ निश्चित करू. सदस्यत्वाच्या व्याख्येनुसार $a \in \fregeext{x}{F(x)}$ तेव्हाच आणि तेव्हाच जेव्हा $\exists G(\fregeext{x}{F(x)}
= \fregeext{x}{G(x)} \land Ga)$; मूलनियम पाचनुसार तेव्हाच आणि तेव्हाच जेव्हा $\exists G(\forall x(Fx \liff Gx) \land Ga)$; आणि प्राथमिक द्वितीय-क्रम तर्कशास्त्रानुसार तेव्हाच आणि तेव्हाच जेव्हा $Fa$.
\end{proof}

यातून अनौपचारिक गुणधर्माधारित संचरचना जवळजवळ लगेच मिळते:

\begin{lem}[\emph{Grundgesetze} मध्ये.]
$\forall F \exists s \forall a (a \in s \liff Fa)$
\end{lem}

\begin{proof}
$F$ निश्चित करू. मग \olref{lem:Fregeextensions} वरून $\forall a (a \in
\fregeext{x}{F(x)} \liff Fa)$ मिळते. म्हणून अस्तित्ववाचक सामान्यीकरणाने $\exists s\forall a(a \in s \liff Fa)$ मिळते. $F$ स्वैर असल्यामुळे अपेक्षित निष्कर्ष मिळतो.
\end{proof}

$\forall x(Fx \liff x \notin x)$ अशी $F$ ची निवड केल्यास रसेलची विरोधापत्ती निष्पन्न होते.

\end{document}
